\documentclass[11pt,a4paper]{article}
\usepackage{fontspec}
\setmainfont{Times New Roman}
\usepackage[russian]{babel}
\usepackage[a4paper,left=17mm,right=17mm,top=14mm,bottom=14mm]{geometry}
\usepackage{amsmath}
\usepackage{xcolor}
\usepackage{needspace}
\pagestyle{empty}
\setlength{\parindent}{0pt}
\setlength{\parskip}{0pt}
\setlength{\emergencystretch}{2em}
\clubpenalty=10000
\widowpenalty=10000
\newcommand{\mixed}[3]{#1\frac{#2}{#3}}
\newcommand{\answerrow}{%
  \par\vspace{3pt}\noindent\small
  Упрощённое выражение: \fbox{\rule{0pt}{17pt}\hspace{132pt}}%
  \hfill Ответ: \fbox{\rule{0pt}{17pt}\hspace{76pt}}\normalsize\par}
\begin{document}
{\small 6 класс\enspace\textperiodcentered\enspace Модуль 2\enspace\textperiodcentered\enspace Задание 7}\par
\vspace{5pt}
{\fontsize{17}{21}\selectfont Умножение смешанных чисел}\par
\vspace{2pt}
{\fontsize{13}{16}\selectfont Счёт, выражения, уравнения}\par
\vspace{7pt}
Для зачёта: 2 из 3.\par
\vspace{7pt}
{\color{black!25}\hrule}
\vspace{10pt}

1. Запишите всю цепочку вычислений и получите ответ.
В окончательном ответе выделите целую часть и сократите дробную.
\[
  \mixed{2}{1}{4}\cdot0{,}8+\frac56\cdot\mixed{1}{3}{5}-\mixed{1}{3}{10}.
\]
\vspace{-4pt}
\noindent{\color{black!30}\rule{\dimexpr\linewidth-2pt\relax}{0.35pt}}\par\vspace{16pt}
\noindent{\color{black!30}\rule{\dimexpr\linewidth-2pt\relax}{0.35pt}}\par
\vspace{12pt}

2. Упростите выражение, затем найдите его значение.
\[
  \mixed{1}{3}{4}x+\left(\frac5{12}x-\mixed{1}{1}{6}\right),\qquad x=1{,}2.
\]
\answerrow\vspace{16pt}

\Needspace{96pt}
3. Какие из произведений являются натуральными числами?
Выпишите номера всех подходящих вариантов. Вычислять значения произведений не нужно.
\[
  \text{1)}\;\mixed{2}{3}{8}\cdot8
  \qquad\qquad
  \text{2)}\;\mixed{1}{3}{5}\cdot\mixed{2}{1}{4}
  \qquad\qquad
  \text{3)}\;\mixed{3}{1}{3}\cdot0{,}3
\]
\vspace{-4pt}
Номера вариантов: \fbox{\rule{0pt}{17pt}\hspace{110pt}}
\end{document}
